% TOP_PROBLEM: {"id": 2602, "title": "Kourovka Notebook Problem 21.93", "category_id": 17, "category": {"id": 17, "name": "group_theory", "display_name": "Group Theory", "description": "Problems about groups, group actions, representations, and related algebraic structures.", "slug": "group-theory", "order_index": 17, "created_at": "2026-07-31T15:26:25.671Z"}, "status": "open", "problem_number": "KOU-21.93", "set_id": 10}
% FIRST_POSTED: 2026-10-01
% ENTRY_KIND: statement_only
% UnsolvedMath ID 2602; source catalogue code KOU-21.93
% !TeX program = lualatex
% Definition and sources only; this file is not an LLM solution attempt.
\documentclass[11pt,a4paper]{article}
\usepackage[margin=25mm]{geometry}
\usepackage{fontspec}
\setmainfont{lmroman10-regular.otf}[BoldFont=lmroman10-bold.otf,
 ItalicFont=lmroman10-italic.otf,BoldItalicFont=lmroman10-bolditalic.otf]
\usepackage{amsmath,amssymb,mathtools}
\usepackage{unicode-math}
\setmathfont{latinmodern-math.otf}
\AtBeginDocument{\let\setminus\smallsetminus}
\usepackage{xurl,hyperref}
\hypersetup{colorlinks=true,urlcolor=blue}
\setcounter{secnumdepth}{0}
\setlength{\parindent}{0pt}
\setlength{\parskip}{5pt}
\setlength{\emergencystretch}{3em}
\begin{document}
\section*{Kourovka Notebook Problem 21.93}\label{2602}
{\small UnsolvedMath ID 2602 · KOU-21.93 · Group Theory · Kourovka Notebook - New Problems, Issue 21}
\subsection{Definitions and mathematical statement}
Let $G$ be a group, with no finiteness assumption, and let $k\ge2$ be an integer. For each $i\in\{1,\ldots,k\}$ choose a subgroup $H_i\le G$ and an element $g_i\in G$. Its left coset is $g_iH_i=\{g_ih:h\in H_i\}$. Suppose these cosets form a partition:
\[
 G=\bigcup_{i=1}^{k}g_iH_i,
 \qquad g_iH_i\cap g_jH_j=\varnothing\quad(i\ne j).
\]
The index $[G:H]$ of a subgroup is the cardinality of the set of its distinct left cosets. The Herzog--Sch\"onheim question asks whether every such partition necessarily has
\[
 [G:H_i]=[G:H_j]\quad\text{for some distinct }i,j.
\]
Thus a finite disjoint covering of a group by at least two subgroup cosets should never have all subgroup indices different.

The partition condition includes disjointness as well as coverage. The subgroups themselves need not be distinct, conjugate, or normal, and the translating elements $g_i$ may differ arbitrarily. Since each coset is nonempty and $k\ge2$, no part can be all of $G$. In the special case of finite $G$, each index equals $|G|/|H_i|$, so the conclusion says that two of the subgroups have equal orders. The formulation with indices also makes sense for infinite ambient groups.
\subsection{Short English statement}
Must every finite partition of a group into at least two subgroup cosets contain two parts whose subgroups have the same index?
\subsection{Sources}
\begin{itemize}
\item[S1] ULAM.AI and the UnsolvedMath Contributors, supplied problems.json, record 2602 (KOU-21.93), Kourovka Notebook - New Problems, Issue 21. Catalogue identity and source transcription; CC BY 4.0.
\url{https://huggingface.co/datasets/ulamai/UnsolvedMath}.
\item[S2] The Kourovka Notebook, 21st edition, version 47 (30 September 2026), new problems, Problem 21.93. Expanded formulation of the supplied catalogue statement.
\url{https://arxiv.org/pdf/1401.0300v47}.
\end{itemize}
\end{document}
